% \newpage
\section*{Задание 7: Миграция кода между технологическими стеками}
\subsection{Постановка задачи}
\noindent
Исходный язык: Python \\
Целевой язык: C++\\
Алгоритм: прямой проход нейронной сети (Forward Propagation)\\
Исходный код: OptimizedNeuralNet из файла v2.py (задание 6)\\
Документация: файл 5.md (README с теоретическим обоснованием)\\
Требование: портировать на C++ с сохранением логики, идиоматичностью и прохождением тестов.\\
\newpage
\subsection{Журнал диалога и итераций}
\noindent
\textbf{Итерация 1: Подготовка контекста и выбор стратегии миграции} \\
\indent
Промпт: 
\textit{
Ты — Senior C++ Engineer, специалист по миграции алгоритмов машинного обучения из Python в высокопроизводительный C++.
Я предоставлю тебе:
1. Исходный код на Python (класс OptimizedNeuralNet с forward pass).
2. Документацию (README.md) с теоретическим обоснованием и примером использования.
Твоя задача — проанализировать исходный код и документацию, выявить ключевые особенности 
(управление памятью, динамическую типизацию, идиомы Python) и предложить стратегию миграции на C++ (C++17/20).
Укажи, какие конструкции Python потребуют замены на идиоматичные C++ (например, списки списков → std::vector,
zip → индексы или диапазоны, транспонирование весов при инициализации). Оцени потенциальные сложности: передача
матриц, производительность, исключения. Выдай результат в виде плана миграции.
} \\
\newline
Ответ модели (Gemini-3):
\begin{figure}[H]
	\centering
	\includegraphics[width=15cm]{z7/1-1.png}
	\caption{Итерация 1. ч.1}
\end{figure}
\newpage
\begin{figure}[H]
	\centering
	\includegraphics[width=17cm]{z7/1-2.png}
	\caption{Итерация 1. ч.2}
\end{figure}
\newpage
\begin{figure}[H]
	\centering
	\includegraphics[width=17cm]{z7/1-3.png}
	\caption{Итерация 1. ч.3}
\end{figure}
\newpage
\begin{figure}[H]
	\centering
	\includegraphics[width=17cm]{z7/1-4.png}
	\caption{Итерация 1. ч.4}
\end{figure}
Анализ результата:
Модель правильно выделила ключевые моменты и предложила разумную стратегию. Нет галлюцинаций. План принимается.
\newpage
\noindent
\textbf{Итерация 2: Идиоматическая трансляция кода} \\
\indent
Промпт:
\textit{
Отлично. Теперь, следуя плану, сгенерируй полный код на C++17 для класса NeuralNetCpp.
Требования:
\begin{itemize}
	\item Используй std::vector, никаких сырых указателей.
	\item Реализуй транспонирование весов в конструкторе так же, как в Python.
	\item Метод forward должен быть константным (не меняет веса).
	\item Для ReLU используй std::max(0.0f, x).
	\item Для Softmax реализуй стабильную версию (вычитание максимума).
	\item Добавь проверку размерностей в конструкторе и в forward.
	\item Код должен быть самодокументированным (комментарии на русском).
	\item Выдай один файл NeuralNetCpp.h с объявлением и определением (header-only для простоты).
\end{itemize}
}
Ответ модели (Gemini-3):
\lstinputlisting[language=C++]{z7/NeuralNetCpp-non-utf.h}
Анализ результата:
Код компилируется, логика полностью соответствует Python-версии. Использованы std::vector, std::max, std::exp. Проверки размерностей присутствуют. Стиль — современный C++. Однако отсутствует тестирование. Переходим к адаптации тестов.]
\newpage
\noindent
\textbf{Итерация 3: Адаптация тестов и верификация} \\
\indent
Промпт:
\textit{
Теперь переведи unit-тесты из Python на C++ с использованием Google Test или простого self-check main. Тестовые сценарии:\\
1. Проверка на примере из документации: вход [[1.0, 0.5]], веса как в README, ожидаемый вывод ~[0.4588, 0.5412]. \\
2. Проверка пустого батча (должно выбросить исключение).\\
3. Проверка неверной размерности входа.\\
4. Проверка на другом случайном наборе (сравнить с результатом Python).\\
Напиши отдельный файл test\_neuralnet.cpp, который подключает NeuralNetCpp.h, запускает тесты и выводит "PASS"\ или "FAIL".
} \\
\newline
Ответ gemini-3 (код тестов):
\lstinputlisting[language=C++]{z7/test_neuralnet-non-utf.cpp}
\noindent
\textbf{Как скомпилировать и запустить:} \\
\indent
Поскольку это header-only решение, вам достаточно положить NeuralNetCpp.h и test\_neuralnet.cpp в одну папку и выполнить:
\begin{verbatim}
					g++ -O3 -std=c++17 test_neuralnet.cpp -o test_nn
					./test_nn
\end{verbatim}
Анализ результата:
При запуске все тесты прошли. Код проходит проверки на срабатывании исключений, также возвращает корректные значения из примера.
\newpage
\noindent
\textbf{Итерация 4: Сравнительный анализ реализаций} \\
\indent
Промпт:
\textit{
Теперь, когда у нас есть работающий C++ код и тесты, выполни сравнительный анализ двух реализаций (Python vs C++). Ответь на вопросы: \\
1. Какие ключевые различия в стиле и идиомах между Python и C++ проявились при портировании?\\
2. Как изменилось управление памятью? Есть ли в C++ утечки?\\
3. Какой код более читаемый?\\
4. Какие оптимизации можно добавить в C++ версию (например, использование SIMD, многопоточность для батча)?\\
5. Оцени, насколько идентична логика (на примере входов и выходов).\\
Выдай структурированный ответ.\\
}\\
\newline
Ответ модели (Gemini-3):
\begin{figure}[H]
	\centering
	\includegraphics[width=15cm]{z7/4-1.png}
	\caption{Итерация 4. ч.1}
\end{figure}
\newpage
\begin{figure}[H]
	\centering
	\includegraphics[width=17cm]{z7/4-2.png}
	\caption{Итерация 4. ч.2}
\end{figure}
\newpage
\begin{figure}[H]
	\centering
	\includegraphics[width=17cm]{z7/4-3.png}
	\caption{Итерация 4. ч.3}
\end{figure}
\noindent
Анализ: Модель дала объективную оценку.
\newpage
\noindent
\textbf{Итерация 5: Формирование финального отчёта о миграции} \\
\indent
Промпт:
\textit{
Собери итоговый отчёт по миграции, включив:
\begin{itemize}
	\item Исходный Python-код (класс OptimizedNeuralNet).
	\item Финальный C++ код (NeuralNetCpp.h).
	\item Лог прохождения тестов (скриншот или текстовый вывод).
	\item Сравнительный анализ (из предыдущего ответа).
	\item Вывод: удалось ли достичь полной функциональной идентичности.
\end{itemize}
}
Ответ Gemini-3 (отчёт):
\lstinputlisting[language={}]{z7/result-non-utf.md}
\subsection{Артефакты}
Финальный C++ код:
\lstinputlisting[language=C++]{z7/NeuralNetCpp-non-utf.h}
Лог прохождения тестов:
\begin{verbatim}
					--- Starting NeuralNetCpp Tests ---
					[PASS] Documentation Example Match
					[PASS] Empty Batch Exception Handled
					[PASS] Dimension Mismatch Exception Handled
					[PASS] ReLU Zeroing and Softmax Unity
					-----------------------------------
					OVERALL STATUS: ALL TESTS PASSED
\end{verbatim}
\subsection{Сравнение моделей}
Финальный промпт (объединяющий все этапы) был подан контрольной модели DeepSeek R1 без изменений.
Результаты сравнения:
\begin{table}[h!]
	\centering
	\begin{tabular}{|l|p{5cm}|p{5cm}|}
		\hline
		\textbf{Критерий} & \textbf{модель (a)} & \textbf{модель (b)} \\
		\hline
		Идентичность логики & 5 & 4 \\
		\hline
		Обработка ошибок & 5 & 5 \\
		\hline
		Качество тестов  & 5 & 4\\
		\hline
		Галлюцинации & нет & нет \\
		\hline
		\textbf{Итог} & 5 & 4.5\\
		\hline
	\end{tabular}
\end{table} \\
gemini-3 более точно следует современным стандартам C++ и даёт более чистый код без лишних деталей. deepseek R1 также справляется, но код иногда выглядит как «калька с Python».
\subsection{Вывод}
Наиболее эффективной техникой промпт-инжиниринга оказалась цепочка рассуждений (Chain-of-Thought) с предварительным анализом различий языков. Также важным было обогащение контекста (передача документации из задания 6). Модель gemini-3 успешно справилась с портированием, сохранив логику и добавив идиоматичные конструкции C++. Контрольная модель deepseek R1 дала рабочий, но менее элегантный код, что подтверждает необходимость выбора основной модели с учётом специфики задачи.